\chapter{\IfLanguageName{dutch}{Requirements verzamelen}{Requirements gathering}}%
\label{ch:requirementsanalyse}

De voorgaande literatuurstudie legde de theoretische fundamenten van\\\gls{esg}-rapportage en \gls{rag} uit. Dit hoofdstuk vertaalt die inzichten naar concrete eisen waaraan de te bouwen dynamische \gls{rag}-pipeline moet voldoen. De requirementsanalyse vormt de schakel tussen de probleemstelling en de technische implementatie: zij bepaalt welke criteria de optimale configuratie moet behalen en biedt een objectief raamwerk voor de latere fases in de vergelijkende studie.

De requirements werden verzameld via twee complementaire bronnen. Enerzijds werden semigestructureerde interviews en gesprekken via e-mail gevoerd met mijn copromotor en opdrachtgever, Gianluca Gueli van Turtle Srl, die als primaire belanghebbende de bedrijfsnoden vertolkt. Anderzijds werd de literatuurstudie aangewend om technische en wetenschappelijke eisen te concretiseren, in het bijzonder op het vlak van \gls{rag}-evaluatie \autocite{Es2024} en architecturale configureerbaarheid. De gecombineerde input werd gestructureerd via de MoSCoW-methode, een prioriteringstechniek die vereisten indeelt in vier categorieën: Must Have, Should Have, Could Have en Won't Have.

\section{Functionele requirements}

Functionele requirements beschrijven wat het systeem concreet moet kunnen doen. Zij zijn direct afgeleid van de bedrijfsnoden van Turtle Srl en de kerntaken van een \gls{rag}-pipeline.

\begin{itemize}
    \item \textbf{Informatie-extractie uit complexe bronnen.} Het systeem moet in staat zijn om relevante data te extraheren uit diverse en ongestructureerde bestandsformaten. \gls{esg}-documenten bevatten informatie die wordt weergegeven in tabellen en tekstuele bijlagen binnen \gls{pdf}- en Excel-bestanden. De \gls{rag}-pipeline moet deze data effectief kunnen parsen en verwerken om inhoudelijk juiste antwoorden te garanderen.
    \item \textbf{Meertalige ondersteuning (IT/EN).} Embeddingmodellen en \glspl{llm} moeten meertalig zijn om diverse \gls{esg}-documenten te kunnen verwerken en antwoorden te genereren.
    \item \textbf{Granulaire bronverwijzing (paginaniveau).} Naast het identificeren van het brondocument is het wenselijk dat het systeem de specifieke pagina vermeldt waar de informatie is gevonden. Dit verhoogt de efficiëntie voor de auditor aanzienlijk bij de handmatige verificatie van de gegenereerde output, aangezien het de zoektijd in omvangrijke rapporten reduceert.
    \item \textbf{Granulaire bronverwijzing op alineaniveau.} Hoewel een verwijzing op alineaniveau de hoogste precisie biedt, valt dit buiten de huidige scope van de \gls{poc}. De technische complexiteit bij het accuraat mappen van tekstfragmenten naar specifieke alinea's in ongestructureerde \gls{pdf}-documenten is hoog en de implementatie hiervan wordt in deze fase niet als kritiek beschouwd voor het aantonen van de basisfunctionaliteit.
\end{itemize}

\section{Niet-functionele requirements}

Niet-functionele requirements beschrijven de kwaliteitseisen en randvoorwaarden waaraan het systeem moet voldoen. Zij zijn bepalend voor de bruikbaarheid en wetenschappelijke geldigheid van het onderzoek.

\begin{itemize}
  \item \textbf{Juridische traceerbaarheid \& auditeerbaarheid.} Bij elk gegenereerd antwoord moet het systeem expliciet aangeven welk brondocument als basis dient. Deze eis werd door mijn copromotor omschreven als essentieel: \textit{``It is essential for the system to clearly indicate which documents it used to generate each answer''}. Vanuit de literatuur wordt deze eis bevestigd door de juridische traceerbaarheidsvereisten van de \gls{csrd}, waarbij elke bewering in een \gls{esg}-rapport herleidbaar moet zijn naar brondocumentatie \autocite{Centre.2023}.
  \item \textbf{Feitelijke betrouwbaarheid (faithfulness).} Antwoorden moeten feitelijk consistent zijn met de opgehaalde brondocumenten om hallucinaties te voorkomen. Deze vereiste wordt geëvalueerd via de faithfulness-metriek van het \gls{ragas}-framework \autocite{Es2024}.
  \item \textbf{Data soevereiniteit \& privacy.} Gezien de vertrouwelijkheid van de verwerkte bedrijfsgegevens is dataverwerking conform de Europese privacywetgeving een strikte vereiste. Om aan de \gls{gdpr}-richtlijnen te voldoen, moet worden gegarandeerd dat data uitsluitend binnen de \gls{eu} wordt opgeslagen en verwerkt. Deze eis stuurt de selectie naar in de \gls{eu} gevestigde providers en de inzet van modellen die lokaal of binnen een private cloudomgeving gehost kunnen worden.
  \item \textbf{Azure-infrastructuur compatibiliteit \& integreerbaarheid.} De \gls{poc} van Turtle Srl wordt gehost op Azure. Dit stelt eisen aan de keuze van modellen en diensten: proprietaire modellen (OpenAI, Google) zijn beschikbaar via \glspl{api} die vanuit Azure aanroepbaar zijn, terwijl open-source modellen een Azure \gls{vm}- of containerdeployment vereisen.
  \item \textbf{Predefined pipelines (reproduceerbaarheid).} Voor de wetenschappelijke integriteit van de thesis moet elke experimentele run gebruikmaken van vooraf gedefinieerde pipelines, zodat resultaten verifieerbaar en herhaalbaar zijn. Dit onderscheidt de experimentele opzet van een real-world productieomgeving, waar een dynamisch configuratie framework zoals HyPSTER de voorkeur verdient.
\end{itemize}

\section{Prioritering via MoSCoW}

Op basis van de bovenstaande requirements werd een geprioriteerde lijst opgesteld via de MoSCoW-methode. Tabel~\ref{tab:moscow} geeft een overzicht.

\begin{table}[h]
  \centering
  \begin{tabular}{p{0.04\textwidth} p{0.06\textwidth} p{0.33\textwidth} p{0.48\textwidth}}
    \toprule
    \textbf{P} & \textbf{Type} & \textbf{Requirement} & \textbf{Verantwoording} \\
    \midrule
    M & NFR & Juridische traceerbaarheid                   & Elke bewering moet direct herleidbaar zijn naar brondocumentatie (essentieel voor \gls{csrd}-compliance). \\
    M & NFR & Feitelijke betrouwbaarheid (faithfulness)    & Hallucinaties zijn onacceptabel in een audit-omgeving; de output moet 100\% steunen op de context. \\
    M & NFR & Data soevereiniteit \& privacy               & Klantdata is strikt vertrouwelijk en moet binnen de \gls{eu}-jurisdictie verwerkt worden (\gls{gdpr}). \\
    M & FR  & Informatie-extractie uit complexe bronnen    & Het vermogen om bruikbare antwoorden te destilleren uit heterogene \gls{pdf}- en Excel-bestanden. \\
    \midrule
    S & FR  & Meertalige ondersteuning (IT/EN)             & Directe verwerking van zowel Italiaanse als Engelse documenten zonder externe vertaalstap. \\
    S & NFR & Azure-infrastructuur compatibiliteit         & De oplossing moet integreerbaar zijn in de bestaande cloud-omgeving van Turtle Srl. \\
    S & NFR & Predefined pipelines (reproduceerbaarheid)   & Wetenschappelijke integriteit van de thesis waarborgen. \\
    \midrule
    C & FR  & Granulaire bronverwijzing                    & Verhoogt het gebruiksgemak voor de auditor maar is geen harde blokkade voor de \gls{poc}. De focus ligt op paginaniveau. \\
    \midrule
    W & FR  & Granulaire bronverwijzing op alineaniveau    & De technische complexiteit van exacte alinea-mapping weegt in deze fase niet op tegen de meerwaarde. \\
    \bottomrule
  \end{tabular}
  \caption[MoSCoW-prioritering van de requirements]{\label{tab:moscow}Geprioriteerde lijst betreffende de requirements voor dit onderzoek.}
\end{table}

De MoSCoW-prioritering maakt duidelijk dat de kern van dit onderzoek rust op vier fundamentele pijlers: \textbf{juridische traceerbaarheid}, \textbf{feitelijke betrouwbaarheid} (\textit{faithfulness}), strikte \textbf{datasoevereiniteit} binnen de \gls{eu}-jurisdictie en de effectieve \textbf{informatie-extractie} uit complexe, heterogene bronnen. Deze vier Must-Have-requirements vormen samen het dwingende kader waartegen alle 27 configuraties in de experimentele fase worden geëvalueerd. Hierbij wordt bepaald welke technische invulling de optimale balans biedt tussen audit-waardigheid, privacy en performantie voor de specifieke casus van Turtle Srl.